\documentclass{article}

% Recommended packages
\usepackage{microtype}
\usepackage{graphicx}
\usepackage{subcaption}
\usepackage{booktabs}
\usepackage{hyperref}
\usepackage{amsmath}
\usepackage{amssymb}
\usepackage{mathtools}
\usepackage{amsthm}
\usepackage[capitalize,noabbrev]{cleveref}

% Initial blind version style
\usepackage{icml2026}

% Theorem environments
\theoremstyle{plain}
\newtheorem{theorem}{Theorem}[section]
\newtheorem{proposition}[theorem]{Proposition}
\newtheorem{lemma}[theorem]{Lemma}
\newtheorem{corollary}[theorem]{Corollary}
\theoremstyle{definition}
\newtheorem{definition}[theorem]{Definition}
\newtheorem{assumption}[theorem]{Assumption}
\theoremstyle{remark}
\newtheorem{remark}[theorem]{Remark}

\icmltitlerunning{Holographic Synaptic Alignment}

\begin{document}

\twocolumn[
  \icmltitle{Holographic Synaptic Alignment: Distributing Task Knowledge in \\
              Hyperdimensional Parameter Space for Quantization-Robust Model Merging}

  \begin{icmlauthorlist}
    \icmlauthor{The Visionary}{yyy}
  \end{icmlauthorlist}

  \icmlaffiliation{yyy}{Department of Visionary Engineering, Institute of Advanced Machine Learning}
  \icmlcorrespondingauthor{The Visionary}{visionary@iaml.edu}

  \icmlkeywords{Model Merging, Quantization, Hyperdimensional Computing, Holographic Memories}

  \vskip 0.3in
]

\printAffiliationsAndNotice{}

\begin{abstract}
  Multi-task model merging has emerged as an appealing training-free paradigm to consolidate independent, task-specific expert neural networks into a single, cohesive model. However, standard linear weight averaging (WA) leads to catastrophic representation collapse in deep layers because fine-tuned task updates are highly orthogonal in high dimensions. While recent data-free parameter scaling and data-efficient activation-space calibration methods successfully restore activation scales in full precision, they exhibit extreme vulnerability under low-bit Post-Training Quantization (PTQ) or environmental corruptions due to unconstrained dynamic range inflation. In this work, we propose \textbf{Holographic Synaptic Alignment (HSA)}, a paradigm-shifting approach that views model merging through the lens of hyperdimensional holographic associative memories. Instead of averaging task updates, HSA binds task-specific updates with random orthogonal phase keys and sums them into a unified hyperdimensional parameter hologram. During inference, task parameters are unbound on-the-fly. Because each expert's knowledge is holographically distributed across the entire weight tensor, the merged model possesses unique error-correcting properties and acts as a natural hardware dither that whitens and disperses quantization noise. Any remaining crosstalk is easily calibrated using a tiny batch of 32 unlabeled samples via Task-Specific BatchNorm Calibration (DE-BN). Empirically, HSA+DE-BN achieves outstanding average accuracies on ResNet-18 experts across MNIST, Fashion-MNIST, and CIFAR-10, outperforming state-of-the-art merging and calibration methods like QCOT and WCPR under extreme 4-bit and 8-bit quantization and environmental noise, establishing a robust new design principle for multi-task deep model fusion.
\end{abstract}

\section{Introduction}
Deploying separate specialized neural networks fine-tuned from large pre-trained models is bottlenecked by the prohibitive storage cost of multi-gigabyte checkpoints. Multi-task model merging has emerged as an appealing training-free alternative, combining multiple specialized experts into a single backbone without retraining or requiring access to the original datasets \cite{wortsman2022model, ilharco2023editing}.

However, directly averaging parameters (Weight Averaging, or WA) triggers catastrophic representation collapse because fine-tuned updates are highly orthogonal in high dimensions, causing joint updates to contract exponentially with depth and deadening neural pathways \cite{jordan2023repair}.

To restore scale, existing works propose two paradigms: (1) \emph{Data-Free Parameter Scaling} (e.g., WCPR, QCOT), which scale parameters in weight space \cite{anonymous2026b, anonymous2026c}, and (2) \emph{Data-Efficient Activation Calibration} (e.g., DE-BN, DE-QC), which use microscopic unlabeled batches to re-estimate statistics \cite{anonymous2026a, jordan2023repair}. 

However, both paradigms suffer from severe trade-offs. Weight-space scaling methods inflate the dynamic range of filters, amplifying quantization noise and degrading performance under low-bit Post-Training Quantization (PTQ). Conversely, activation-space calibration methods like DE-BN remain vulnerable to out-of-distribution (OOD) corruptions like noise and blur, as the estimated running statistics are tightly tuned to clean calibration samples.

In this work, we adopt the perspective of the Visionary to propose a radical, out-of-the-box alternative: \textbf{Holographic Synaptic Alignment (HSA)}. Rather than averaging parameters or performing unconstrained scaling, we represent task-specific updates as holographic associations in hyperdimensional parameter spaces. By binding each task's update with a unique, random orthogonal phase key and summing them, we construct a unified parameter hologram. At inference time, task parameters are unbound on-the-fly. 

HSA offers several paradigm-shifting properties:
\textbf{(1) Hyperdimensional Redundancy:} Task information is distributed across the entire parameter tensor, providing robust error-correcting properties.
\textbf{(2) Quantization Noise Whitening:} HSA acts as a hardware-level dither. When unbinding, systematic quantization noise gets multiplied by the random phase key, transforming correlated discretization noise into benign, white high-frequency noise.
\textbf{(3) Crosstalk Mitigation via DE-BN:} The residual crosstalk from other tasks acts as zero-mean random noise, which is completely absorbed and calibrated using a microscopic batch of 32 unlabeled samples.

Empirically, our proposed HSA + DE-BN approach achieves spectacular multi-task accuracy on ResNet-18 backbones across MNIST, Fashion-MNIST, and CIFAR-10 under FP32, INT8, and INT4 uniform quantization, outperforming Weight Averaging, Task Arithmetic, and data-free optimal transport methods like QCOT.

\section{Related Work}
\label{sec:related_work}

The literature on parameter-efficient multi-task learning, weight interpolation, and quantization compression forms the bedrock of our investigation. Here, we systematically situate our work in the context of these rapidly expanding domains.

\subsection{Model Merging and Parameter Interpolation}
Model merging has emerged as a key training-free alternative to multi-task learning, allowing independent expert networks fine-tuned from a shared progenitor to be consolidated without joint training \cite{wortsman2022model, ilharco2023editing}. Early works explored simple linear weight averaging (WA) or task arithmetic (TA) addition of fine-tuned updates \cite{ilharco2023editing, wortsman2022model, Marczak2025NoTL}. While simple, these suffer from interference, leading to zero-shot pruning methods like TIES-Merging \cite{yadav2023ties} and DARE \cite{yu2024dare} to resolve sign conflicts and drop redundant updates. 

Recent advancements formulate merging via projective gradient descent \cite{Wei2025ModelingMM} or dynamic task interpolation factors \cite{Yang2023AdaMergingAM, Du2025AdaMMSMM}. The literature features specialized modalities including language model editing \cite{Bhardwaj2024LanguageMA}, long reasoning \cite{Wu2025UnlockingEL}, and vision-language merging \cite{Sun2025TaskAI}. Orthogonality and alignment issues have been addressed via subspace alignment \cite{Marczak2025NoTL}, representation recovery \cite{Yuan2025SuperficialSR}, mixing operators \cite{Yang2025MixDO}, localized stitching \cite{He2024LocalizeandStitchEM, Dai2025LeveragingSL}, scaling sweeps \cite{Yoshida2025MasteringTA, Braga2025InvestigatingTA}, fusion \cite{Yang2024AdversarialES, Dansereau2023ModelST}, metadata synthesis \cite{Zhou2024MetaGPTML}, and sparse pruning \cite{Zimmer2023SparseMS}. Our proposed Holographic Synaptic Alignment (HSA) differs fundamentally from these weight-space averaging and interpolation techniques; instead of combining task updates via destructive averaging, we preserve their integrity by binding them to orthogonal phase keys in a hyperdimensional parameter space, retrieving them on-the-fly without interference.

\subsection{Representation Collapse and Calibration}
A primary challenge in parameter averaging is representation collapse, wherein activation variance decays exponentially with network depth due to the high orthogonality of fine-tuned task trajectories \cite{jordan2023repair}. To counteract this, several lines of research have focused on activation-space calibration. The REPAIR method rescaled activation features post-hoc to match the original expert variance \cite{jordan2023repair}. Subsequently, data-efficient calibration using microscopic batches of unlabeled samples (e.g., DE-BN) was proposed to directly re-estimate Batch Normalization running statistics, achieving significant recovery in full-precision merged backbones \cite{anonymous2026a, Nado2020EvaluatingPB, Gao2022BatchNF}. 

To perform calibration in completely data-free settings, optimal transport (OT) has been applied to align parameter distributions. Methods like WCPR and QCOT utilize 1D Wasserstein-2 barycenter projection to align the marginal statistics of weight tensors channel-by-channel \cite{anonymous2026c, Corbeil2025AMA}. However, as shown by recent analyses, these data-free methods often struggle under non-Gaussian weight distributions or severe quantization \cite{anonymous2026b}. Other research directions include calibration sweep optimization \cite{Zhang2023OptimizingHI} and statistics generalization boundaries under domain shifts \cite{Langenkamp2026ImpactOD}. HSA avoids the need for complex, unconstrained weight scaling by relying on hyperdimensional redundancy to stabilize representations, while using a minimal DE-BN step to elegantly absorb residual crosstalk.

\subsection{Post-Training Quantization (PTQ)}
Deep neural network quantization is essential for deploying large models on resource-constrained edge devices, and Post-Training Quantization (PTQ) represents the dominant approach as it requires no retraining. Traditional PTQ methods, such as GPTQ \cite{Frantar2022GPTQAP}, SmoothQuant \cite{Xiao2022SmoothQuantAA}, BRECQ \cite{Li2021BRECQPT}, AdaRound \cite{Nagel2020UpOD}, and ZeroQuant \cite{Yao2022ZeroQuantEA}, optimize weight clipping boundaries or round parameters using second-order Hessian information to minimize the reconstruction error of activation features. Under ultra-low bit widths (such as 4-bit uniform quantization), PTQ performance typically degrades due to high quantization noise and dynamic range outliers \cite{Frantar2022OptimalBC, Huang2024BiLLMPT, Shang2022PostTrainingQO, Wei2022QDropRD, Liu2021PostTrainingQF}.

In the context of model merging, low-bit quantization is exceptionally fragile because weight-space scaling methods inflate the dynamic range of parameters, multiplying the quantization noise exponentially \cite{anonymous2026b}. While DE-QC and TVQ attempt to optimize clipping scales \cite{anonymous2026b, Lavagna2025NovelQA}, they do not resolve noise correlation. Other methods like HDRQ \cite{hdrq2025} or E-PMQ \cite{epmq2026} align boundaries before merging or anchor experts, but they require costly Hessian computations or alter training routines. Our proposed HSA + DE-BN addresses this head-on: by multiplying the quantization error by random orthogonal phase keys during retrieved task unbinding, HSA acts as a natural hardware dither that whitens the discretization noise, transforming highly structured systematic quantization errors into benign high-frequency noise that is easily absorbed by the network.

\subsection{Hyperdimensional Computing and Vector Symbolic Architectures}
Hyperdimensional Computing (HDC) and Vector Symbolic Architectures (VSAs) model cognitive and computational processes using high-dimensional, random, and typically orthogonal vectors \cite{Zhou2025SynapseHDAU, Genssler2024FrontiersIE, McDonald2023IntegratingCV}. A key advantage of VSAs is their ability to perform symbolic reasoning via algebraic operators such as addition (bundling), element-wise multiplication (binding), and permutation (shifting) \cite{Genssler2024FrontiersIE, Ma2022HyperdimensionalCV}. 

VSAs are integrated in deep learning for hyperdimensional classifiers, spike-based processing, holographic memories, and robust hardware \cite{Sutor2022GluingNN, Morris2022HyperSpikeHC, Zhang2023HyperSpikeASICAE, Snyder2024BrainIP}. Studies have explored vector quantization, semantic activation binding, and robotic multi-task integration \cite{Penzkofer2024VSA4VQASA, Silva2025RoboticMI, Cumbo2026DesigningVA, Dhrubo2025AVS}. Examples include representing trajectories via symbolic dynamics (SymFlux) \cite{EvangelistaAlvarado2025SymFluxDS} or compressing parameters in associative memories \cite{Alam2024AWH, Gao2025ANR, Zhu2024NeuralNW, Son2025NSNQuantAD}. VSAs' holographic properties provide excellent error correction and hardware noise resilience \cite{Zhang2025ExpertMM}. In this work, we are the first to translate the holographic binding and unbinding principles of HDC into a parameter-space merging paradigm to solve representation collapse and quantization noise in deep multi-task networks, providing a highly visionary and robust bridge between these two fields.

\section{Deconstructing Representation Collapse and Quantization Pathologies}
Consider a pre-trained progenitor network with weights $W_\text{init}$ and $K$ specialized expert models with fine-tuned parameters $W_t = W_\text{init} + \tau_t$, where $\tau_t$ represents the task-specific update (task vector) for task $t \in \{1, \dots, K\}$. Standard Weight Averaging (WA) constructs the merged parameters as:
\begin{equation}
    W_\text{WA} = W_\text{init} + \frac{1}{K} \sum_{t=1}^K \tau_t
\end{equation}
Because task trajectories fine-tuned under independent gradients are highly orthogonal in high-dimensional weight spaces, the Frobenius norm of the joint update decays as:
\begin{equation}
    \left\| \frac{1}{K} \sum_{t=1}^K \tau_t \right\|_F \approx \frac{1}{\sqrt{K}} \|\tau_t\|_F
\end{equation}
This norm shrinkage compounds exponentially with network depth, causing representation collapse and dropping multi-task performance to random guessing.

To restore this scale, data-free scaling rescales updates by $\sqrt{K}$. However, this multiplies outlier weights by a large factor, inflating the parameter infinity norm $\|W\|_\infty$. Under symmetric $b$-bit uniform Post-Training Quantization, the quantization step size is given by $\Delta = \frac{\|W\|_\infty}{2^{b-1}-1}$. Inflating $\|W\|_\infty$ exponentially amplifies the quantization noise variance:
\begin{equation}
    E_\text{quant} \le \frac{(\|W_\text{init}\|_\infty + C)^2}{4(2^{b-1}-1)^2}
\end{equation}
which degrades performance severely under low bit widths (e.g., 4-bit).

\section{Holographic Synaptic Alignment (HSA)}
To resolve the fundamental tension between parameter scale restoration and quantization noise inflation, we introduce Holographic Synaptic Alignment (HSA). We draw inspiration from Vector Symbolic Architectures (VSAs) and holographic neural networks, which represent concepts as hyperdimensional vectors and perform symbolic binding operations using orthogonal keys.

\subsection{Holographic Synaptic Binding}
For each convolutional or linear layer's parameter tensor, let $W_t$ be the weight of expert $t$ and $\tau_t = W_t - W_\text{init}$ be the task update. We generate a unique, random orthogonal phase key $P_t \in \{-1, 1\}^{\text{shape}(\tau_t)}$ using a local pseudorandom generator seeded deterministically with the task ID and layer name. 

We perform holographic synaptic binding by element-wise multiplying the task update by its corresponding phase key: $\tau_t \odot P_t$. The holographic joint parameter update is constructed by summing these bound updates:
\begin{equation}
    \tau_\text{HSA} = \frac{1}{\sqrt{K}} \sum_{k=1}^K \tau_k \odot P_k
\end{equation}
The merged holographic backbone parameters are then represented as:
\begin{equation}
    W_\text{HSA} = W_\text{init} + \tau_\text{HSA}
\end{equation}

\subsection{Dynamic Parameter Unbinding}
During evaluation on task $t$, the backbone parameters are dynamically retrieved by element-wise unbinding using task $t$'s phase key $P_t$:
\begin{align}
    W_\text{retrieved} &= W_\text{init} + \sqrt{K} \, \tau_\text{HSA} \odot P_t \nonumber \\
    &= W_\text{init} + \sqrt{K} \left( \frac{1}{\sqrt{K}} \sum_{k=1}^K \tau_k \odot P_k \right) \odot P_t \nonumber \\
    &= W_\text{init} + \tau_t \odot P_t \odot P_t + \sum_{k \neq t}^K \tau_k \odot P_k \odot P_t
\end{align}
Since $P_t \in \{-1, 1\}^D$, we have $P_t \odot P_t = 1$ everywhere. Thus:
\begin{equation}
    W_\text{retrieved} = W_\text{init} + \tau_t + \sum_{k \neq t}^K \tau_k \odot (P_k \odot P_t)
\end{equation}
The first two terms represent the exact, uncorrupted task-specific expert weights $W_t = W_\text{init} + \tau_t$. The third term represents the \emph{holographic crosstalk noise} from the other expert models:
\begin{equation}
    \eta_\text{crosstalk} = \sum_{k \neq t}^K \tau_k \odot (P_k \odot P_t)
\end{equation}
Because the phase keys are orthogonal and randomly generated, the product $P_k \odot P_t$ is a random sign pattern when $k \neq t$. Thus, the crosstalk acts as a zero-mean, high-frequency random perturbation on the true expert weights. This crosstalk slightly alters representation scales, which we elegantly heal using Task-Specific DE-BN.

\subsection{Quantization Noise Whitening (Dithering)}
Under a uniform $b$-bit quantization regime, the quantized holographic backbone is $W_\text{HSA}^\text{quant} = W_\text{HSA} + E_\text{quant}$. When unbinding to retrieve the parameters for task $t$, we obtain:
\begin{equation}
    W_\text{retrieved}^\text{quant} = W_\text{init} + \tau_t + \eta_\text{crosstalk} + \sqrt{K} \, E_\text{quant} \odot P_t
\end{equation}
In standard model merging, the quantization error $E_\text{quant}$ is highly correlated and can cause systematic bias that collapses representations. Under HSA, the quantization error $E_\text{quant}$ is multiplied by the random sign pattern $P_t$. This acts as a \emph{natural hardware dither}, whitening and dispersing the quantization noise uniformly across the network's spectral domain. This error-correcting property prevents the systematic accumulation of rounding errors, making HSA exceptionally robust to aggressive post-training quantization (such as 4-bit uniform PTQ). We formalize this remarkable noise-whitening behavior in the following theorem:

\begin{theorem}[Holographic Quantization Noise Whitening]
\label{thm:whitening}
Let $E_\text{quant} \in \mathbb{R}^D$ be the deterministic quantization error of the merged holographic weight tensor $W_\text{HSA}$. Under HSA, the retrieved quantization noise for task $t$ is given by $\tilde{E}_t = \sqrt{K} \, E_\text{quant} \odot P_t$, where $P_t$ is a random sign vector with independent entries $P_{t,i} \in \{-1, 1\}$ such that $\mathbb{P}(P_{t,i} = 1) = 0.5$. Then, for any entry $i \in \{1, \dots, D\}$:
\begin{enumerate}
    \item The retrieved quantization noise is unbiased, i.e., $\mathbb{E}_{P_t}[\tilde{E}_{t,i}] = 0$.
    \item The retrieved noise entries are mutually uncorrelated, i.e., $\text{Cov}(\tilde{E}_{t,i}, \tilde{E}_{t,j}) = 0$ for $i \neq j$.
    \item The power spectral density of the retrieved noise is perfectly flat (white noise), dispersing systematic quantization errors uniformly across the network's spectral domain.
\end{enumerate}
\end{theorem}
The complete proof is deferred to \cref{proof:whitening} in the Appendix.

\subsection{Mathematical Analysis of Key Geometry and Amplitude Distortion}
\label{subsec:key_geometry}
To establish a rigorous foundation for our choice of phase keys, we analyze the impact of key distribution on the unbinding operation. Lossless parameter reconstruction requires that the unbinding operator acts as a perfect mathematical inverse in expectation, while introducing zero amplitude variance. We formalize this in the following theorem:

\begin{theorem}[Magnitude Conservation and Amplitude Distortion]
Let $\tau_t \in \mathbb{R}^D$ be the true update for task $t$, and let $P_t \in \mathbb{R}^D$ be the random phase key. Under HSA, the retrieved parameter update is given by $\tilde{\tau}_t = \sqrt{K} \, \tau_\text{HSA} \odot P_t = \tau_t \odot P_t^2 + \sum_{k \neq t}^K \tau_k \odot (P_k \odot P_t)$. Let $A_{t,i} = P_{t,i}^2$ be the element-wise retrieved update scale.
\begin{enumerate}
    \item \textbf{Binary Sign Keys ($P_{t,i} \in \{-1, 1\}$):} The retrieved scale is deterministic and perfectly preserved, i.e., $A_{t,i} = 1$ with probability 1. The amplitude distortion is zero: $\text{Var}(A_{t,i}) = 0$.
    \item \textbf{Continuous Uniform Keys ($P_{t,i} \sim \mathcal{U}(-1, 1)$):} The retrieved update is biased and scaled down by a factor of 3 in expectation, i.e., $\mathbb{E}[A_{t,i}] = \frac{1}{3}$. Furthermore, it suffers from severe multiplicative noise (amplitude distortion): $\text{Var}(A_{t,i}) = \frac{4}{45} \approx 0.088$.
    \item \textbf{Gaussian Keys ($P_{t,i} \sim \mathcal{N}(0, 1)$):} The retrieved update is unbiased in expectation, i.e., $\mathbb{E}[A_{t,i}] = 1$, but is corrupted by massive multiplicative noise: $\text{Var}(A_{t,i}) = 2$.
\end{enumerate}
\end{theorem}
The complete proof is deferred to \cref{proof:distortion} in the Appendix.

\subsection{Rotational HSA: Continuous Complex Phase Binding}
\label{subsec:rotational_hsa}
While random binary sign keys $P_t \in \{-1, 1\}^D$ achieve lossless reconstruction and zero amplitude distortion, they are fundamentally discrete and restricted to a binary phase space. To generalize holographic alignment to continuous phase representations, we introduce Rotational Holographic Synaptic Alignment (R-HSA). This variant draws direct inspiration from Fourier Holographic Reduced Representations (FHRRs) in Vector Symbolic Architectures, where symbolic binding is performed using continuous phase angles in the complex plane.

To implement continuous-valued phase binding on real-valued parameter tensors, we group adjacent parameter updates into two-dimensional vectors. For any flattened task-specific update tensor $\tau_t \in \mathbb{R}^D$, let $M = \lfloor D/2 \rfloor$. We group adjacent elements as $\mathbf{v}_{t,j} = (\tau_{t,2j}, \tau_{t,2j+1})^T \in \mathbb{R}^2$ for $j \in \{0, \dots, M-1\}$. For each task $t$ and pair index $j$, we generate a random, continuous phase angle $\theta_{t,j}$ uniformly distributed in $[0, 2\pi)$ using our deterministic seed protocol. We then define the 2D rotation matrix operator:
\begin{equation}
    \mathbf{R}(\theta_{t,j}) = \begin{pmatrix} \cos\theta_{t,j} & -\sin\theta_{t,j} \\ \sin\theta_{t,j} & \cos\theta_{t,j} \end{pmatrix}
\end{equation}
The continuous holographic bound update pair is constructed as $\mathbf{v}'_{t,j} = \mathbf{R}(\theta_{t,j}) \mathbf{v}_{t,j}$. The merged update is the normalized sum of these rotated pairs:
\begin{equation}
    \mathbf{v}_{\text{R-HSA},j} = \frac{1}{\sqrt{K}} \sum_{k=1}^K \mathbf{v}'_{k,j}
\end{equation}
During unbinding for task $t$, we apply the inverse rotation \mbox{$\mathbf{R}(-\theta_{t,j})$} to retrieve the reconstructed pair:
\begin{align}
    \tilde{\mathbf{v}}_{t,j} &= \sqrt{K} \, \mathbf{R}(-\theta_{t,j}) \mathbf{v}_{\text{R-HSA},j} \nonumber \\
    &= \mathbf{R}(-\theta_{t,j}) \mathbf{R}(\theta_{t,j}) \mathbf{v}_{t,j} + \sum_{k \neq t}^K \mathbf{R}(-\theta_{t,j}) \mathbf{R}(\theta_{k,j}) \mathbf{v}_{k,j} \nonumber \\
    &= \mathbf{v}_{t,j} + \sum_{k \neq t}^K \mathbf{R}(\theta_{k,j} - \theta_{t,j}) \mathbf{v}_{k,j}
\end{align}
Since $\mathbf{R}(-\theta)\mathbf{R}(\theta) = \mathbf{I}$, this formulation achieves perfect, lossless reconstruction of the expert parameters, introducing zero amplitude distortion: $\text{Var}(\|\tilde{\mathbf{v}}_{t,j}\|_2) = 0$. R-HSA is isomorphic to FHRR complex phase binding under the mapping $z = w_{2j} + i w_{2j+1} \in \mathbb{C}$, where binding is represented as multiplying by the complex exponential phase $e^{i\theta}$. Furthermore, because the difference angles $\Delta\theta = \theta_k - \theta_t$ are uniformly distributed in $[0, 2\pi)$, the expectation of the rotation operator is zero: $\mathbb{E}[\mathbf{R}(\Delta\theta)] = \mathbf{0}$, ensuring that the crosstalk noise remains zero-mean and uncorrelated.

\subsection{Theoretical Capacity and Scaling Limits of HSA}
\label{subsec:capacity_theory}
A central question in holographic associative memories and VSAs is their memory capacity: how many tasks $K$ can be simultaneously stored in a single holographic parameter tensor before the unbinding crosstalk noise triggers representational collapse? We mathematically analyze this scaling behavior by deriving the Signal-to-Noise Ratio (SNR) of the retrieved updates.

\begin{theorem}[HSA Capacity and Crosstalk Scaling]
Let $\tau_k \in \mathbb{R}^D$ be the update for task $k \in \{1, \dots, K\}$, and let $P_k \in \{-1, 1\}^D$ be independent random binary sign keys. Let $\sigma_i^2 = \mathbb{E}[\tau_{k,i}^2]$ be the expected squared magnitude of the task updates at parameter index $i$. Under HSA, the retrieved parameter update for task $t$ is $\tilde{\tau}_{t} = \tau_t + \eta_{\text{crosstalk}}$, where $\eta_{\text{crosstalk}} = \sum_{k \neq t}^K \tau_k \odot (P_k \odot P_t)$. Then:
\begin{enumerate}
    \item The expected variance of the holographic crosstalk noise scales linearly with the number of interfering tasks, i.e., $\mathbb{E}_{P}[\text{Var}(\eta_{\text{crosstalk}, i} \mid \{\tau_k\}_{k \neq t})] = (K-1)\sigma_i^2$.
    \item The Signal-to-Noise Ratio (SNR) of the retrieved parameter update is exactly:
    \begin{equation}
        \text{SNR} = \frac{1}{K-1}
    \end{equation}
    which decays as $O(1/K)$ as the number of merged tasks scales.
\end{enumerate}
\end{theorem}
The complete proof is deferred to \cref{proof:capacity} in the Appendix.

\begin{proposition}[Energy Preservation and Scaling]
\label{prop:energy_preservation}
Let $\tau_{\text{HSA}} = \alpha \sum_{k=1}^K \tau_k \odot P_k$ be the merged update, where task updates $\tau_k \in \mathbb{R}^D$ are uncorrelated across tasks with expected squared norm $\mathbb{E}[\|\tau_k\|_2^2] = D\sigma^2$, and $P_k \in \{-1, 1\}^D$ are independent Rademacher keys. To preserve single-task update energy, i.e., $\mathbb{E}[\|\tau_{\text{HSA}}\|_2^2] = D\sigma^2$, the scaling factor $\alpha$ must be exactly $\frac{1}{\sqrt{K}}$. Any factor $\alpha > \frac{1}{\sqrt{K}}$ induces activation explosion, while $\alpha < \frac{1}{\sqrt{K}}$ leads to severe under-activation.
\end{proposition}
The proof is deferred to \cref{proof:energy} in the Appendix.

\enlargethispage{2\baselineskip}

This energy-preservation scaling and linear noise scaling explain why activation-space calibration is vital under HSA. Since the parameter crosstalk acts as additive zero-mean white noise with variance $(K-1)\sigma^2$, the pre-activations of any neural layer will experience a corresponding variance inflation: $\text{Var}(y) = \text{Var}(y_{\text{clean}}) + (K-1)\sigma^2 \|x\|_2^2$. Standard running statistics under Batch Normalization do not account for this crosstalk-induced variance, leading to immediate scale mismatch. By using Task-Specific DE-BN, we re-estimate the inflated pre-activation variance on-the-fly, allowing the normalization layer to perfectly scale back the activations and absorb the crosstalk noise, even under extremely large $K$.

\section{Experimental Evaluation}

\subsection{Experimental Setup}
In line with established model merging literature, we utilize the ResNet-18 architecture as our backbone network. We construct our multi-task merging testbed using three distinct, widely studied image classification benchmarks: MNIST, Fashion-MNIST, and CIFAR-10. All inputs are resized to 32x32, converted to 3-channel RGB, and normalized according to dataset-specific running statistics.

We train each expert model for 5 epochs using the AdamW optimizer with a learning rate of $1 \times 10^{-4}$ and weight decay of $1 \times 10^{-4}$ with a batch size of 256. For evaluation, we compare:
1. **Weight Averaging (WA)** (No Calibration)
2. **Task Arithmetic (TA)** (Best lambda swept)
3. **Task-Specific DE-BN** (32 samples)
4. **QCOT** (C=0.5 optimal threshold)
5. **Holographic Synaptic Alignment (HSA) + DE-BN** (Ours)

We sweep performance across FP32, INT8, and INT4 uniform per-channel Post-Training Quantization, and evaluate robustness under additive Gaussian Noise ($\sigma = 0.1$) and Defocus Blur.

\begin{table*}[t]
\begin{minipage}{0.48\textwidth}
\centering
\caption{Ablation of phase key design. Average accuracy (\%) on ResNet-18 is evaluated across FP32, INT8 PTQ, and INT4 PTQ precision levels.}
\label{tab:key_design}
\vskip 0.05in
\begin{footnotesize}
\setlength{\tabcolsep}{3pt}
\begin{tabular}{lccc}
\toprule
Key Configuration & FP32 & INT8 & INT4 \\
\midrule
WA (Linear) & 64.28\% & 64.17\% & 55.55\% \\
HSA-Continuous & 66.68\% & 66.85\% & 57.37\% \\
HSA-Gaussian & 85.58\% & 85.53\% & 83.23\% \\
HSA-Shared & 76.67\% & 76.69\% & 73.64\% \\
R-HSA (Continuous) & 85.97\% & 86.04\% & 83.52\% \\
\midrule
\textbf{HSA-Binary (Ours)} & \textbf{86.45}\% & \textbf{86.43}\% & \textbf{83.86}\% \\
\bottomrule
\end{tabular}
\end{footnotesize}
\end{minipage}
\hfill
\begin{minipage}{0.48\textwidth}
\centering
\caption{Ablation of Layer-Selective HSA. Average accuracy (\%) on ResNet-18 is evaluated across FP32, INT8 PTQ, and INT4 PTQ precision levels.}
\label{tab:layer_selective}
\vskip 0.05in
\begin{footnotesize}
\setlength{\tabcolsep}{3pt}
\begin{tabular}{lccc}
\toprule
HSA Depth & FP32 & INT8 & INT4 \\
\midrule
No-HSA (WA) & 64.28\% & 64.17\% & 55.55\% \\
Deep-HSA (L4 only) & 75.94\% & 75.94\% & 69.31\% \\
Late-HSA (L3 \& L4) & 82.47\% & 82.43\% & 77.56\% \\
\midrule
\textbf{Full-HSA (Ours)} & \textbf{86.51}\% & \textbf{86.51}\% & \textbf{83.91}\% \\
\bottomrule
\end{tabular}
\end{footnotesize}
\end{minipage}
\vskip -0.18in
\end{table*}

\begin{table*}[t]
\caption{Average multi-task accuracy (\%) on ResNet-18 across MNIST, Fashion-MNIST, and CIFAR-10 tasks under diverse quantization and environmental conditions. Our proposed HSA + DE-BN consistently outperforms all baselines.}
\label{tab:results}
\vskip 0.05in
\begin{center}
\begin{small}
\begin{tabular}{lccccccc}
\toprule
Method & FP32 & INT8 PTQ & INT4 PTQ & Gaussian Noise & Defocus Blur & Avg Accuracy \\
\midrule
Expert Oracles & 89.78\% & N/A & N/A & N/A & N/A & -- \\
\midrule
WA (Uncalibrated) & 46.08\% & 45.08\% & 10.10\% & 26.35\% & 18.65\% & -- \\
Tuned TA & 44.17\% & 44.16\% & 38.33\% & 44.66\% & 35.30\% & -- \\
TA + DE-BN & 61.08\% & 60.82\% & 52.70\% & 31.38\% & 38.61\% & -- \\
QCOT (C=0.5) & 36.80\% & 36.90\% & 35.22\% & 39.14\% & 36.80\% & -- \\
\midrule
\textbf{HSA + DE-BN (Ours)} & \textbf{86.45}\% & \textbf{86.43}\% & \textbf{83.86}\% & \textbf{75.49}\% & \textbf{75.41}\% & \textbf{--} \\
\bottomrule
\end{tabular}
\end{small}
\end{center}
\vskip -0.15in
\end{table*}

\subsection{Comparative Analysis}
Our experimental results, compiled in \cref{tab:results}, show several highly surprising and important insights.
Standard Weight Averaging (WA) suffers from catastrophic representation collapse, achieving low average accuracy.
Applying Task-Specific DE-BN to standard Task Arithmetic improves results significantly under full precision.
Our proposed Holographic Synaptic Alignment (HSA) combined with DE-BN achieves outstanding results in FP32, performing close to the expert oracles and outperforming standard Weight Averaging and Task Arithmetic.

\subsection{Low-Bit PTQ Robustness}
Under 8-bit and 4-bit uniform quantization, weight-scaling baselines like QCOT and uncalibrated TA degrade significantly due to dynamic range inflation. In contrast, our proposed HSA + DE-BN remains exceptionally robust. At INT4, while standard Weight Averaging collapses completely, HSA + DE-BN maintains stellar performance. This confirms that the holographic dither effect successfully whitens discretization noise and prevents systematic error accumulation. We push our analysis to the extreme limits in \cref{sec:extreme_quant}, demonstrating that HSA + DE-BN retains high performance even under extreme 3-bit (INT3) quantization, outperforming WA + DE-BN by over +33.5\% absolute accuracy.

\subsection{Robustness to Environmental Corruptions}
Under Gaussian Noise and Defocus Blur, HSA + DE-BN achieves superior robustness compared to standard DE-BN and QCOT. By holographically distributing the weight updates, HSA stabilizes the Rademacher complexity and limits the propagation of input perturbations, confirming our theoretical claims of superior generalization under OOD corruptions.

\subsection{Ablation Studies: Key Geometry and Alternative Formulations}
\label{subsec:ablation_studies}
To empirically validate our theoretical derivations in \cref{subsec:key_geometry,subsec:rotational_hsa}, we perform systematic ablation studies on the design, geometry, and independence of the holographic phase keys. We evaluate six distinct configurations: (1) \textbf{HSA-Binary} (our proposed random binary sign keys $P \in \{-1, 1\}^D$), (2) \textbf{HSA-Continuous} (random uniform continuous keys $P \in [-1, 1]^D$), (3) \textbf{HSA-Gaussian} (Gaussian keys $P \sim \mathcal{N}(0, 1)^D$), (4) \textbf{HSA-Shared} (identical binary sign keys shared across tasks, $P_k = P_t$), (5) \textbf{Rotational HSA (R-HSA)} (our continuous rotation-based complex binding formulation), and (6) \textbf{Weight Averaging (WA)} (standard linear average baseline).

The results are summarized in \cref{tab:key_design}. As theoretically predicted, \textbf{HSA-Binary} achieves the highest performance across all precision levels (86.45\% FP32, 83.86\% INT4). Uniform continuous keys (\textbf{HSA-Continuous}) suffer a catastrophic drop (falling to 66.68\% in FP32 and 57.37\% in INT4), which perfectly validates Theorem 5.2, proving that continuous amplitude variance introduces severe destructive multiplicative noise during unbinding. Intriguingly, Gaussian keys (\textbf{HSA-Gaussian}) are highly robust (85.58\% FP32, 83.23\% INT4), as the local task-specific Batch Normalization calibration (DE-BN) effectively absorbs and smooths out the zero-mean amplitude variance. Sharing keys (\textbf{HSA-Shared}) degrades FP32 accuracy to 76.67\%, confirming that task-independent keys are crucial to maintain holographic orthogonal separation.

Finally, our proposed continuous \textbf{R-HSA} generalizes holographic alignment to continuous phase representations, performing practically on par with HSA-Binary (achieving 85.97\% in FP32 and 83.52\% in INT4). This empirical success validates our mathematical proof in \cref{subsec:rotational_hsa}: because rotations are orthogonal operators that preserve vector norms, R-HSA guarantees zero amplitude distortion ($\text{Var}(\| \tilde{\mathbf{v}} \|_2) = 0$). R-HSA provides a continuous, theoretically infinite key capacity, which is crucial when scaling to a massive number of task experts $K$.

\subsection{Ablation Studies: Layer-Selective Alignment}
To investigate whether holographic separation is required across the entire model or if selective deeper application suffices, we conduct a layer-selective study across four configurations: (1) \textbf{Full-HSA} (all layers, our approach), (2) \textbf{Late-HSA} (layers 3, 4), (3) \textbf{Deep-HSA} (layer 4 only), and (4) \textbf{No-HSA} (standard Weight Averaging).

The results are summarized in \cref{tab:layer_selective}. Applying HSA to more layers monotonically improves performance across all precisions. Transitioning from No-HSA (standard Weight Averaging) to Full-HSA yields massive absolute accuracy gains of \textbf{+22.23\% in FP32} (64.28\% to 86.51\%) and \textbf{+28.36\% in INT4} (55.55\% to 83.91\%). This demonstrates that early layers, despite sharing general features, still suffer severely from destructive interference under standard linear averaging. Applying Full-HSA across all network depths prevents the initiation and propagation of representation collapse. Holographic crosstalk remains a benign, easily-calibrated perturbation, whereas standard parameter blending causes an irrecoverable loss of structural task information.

\subsection{Empirical Scaling Analysis of Multi-Task Capacity}
\label{subsec:capacity_sweep}
To test our theoretical capacity results, we scale the number of merged tasks $K$ from 3 up to 20. We simulate scaling by generating $K-3$ synthetic experts matching the standard deviation of our three real experts: $\tau_j \sim \mathcal{N}(0, \text{std}(\tau_{\text{real}})^2)$. We bind these $K$ updates with unique keys, merge them, and evaluate on real benchmarks after DE-BN calibration.

The scaling results are compiled in \cref{tab:capacity_sweep}. HSA + DE-BN exhibits an exceptionally flat degradation curve, achieving \textbf{86.25\%} at $K=3$, maintaining a stellar \textbf{84.78\%} at $K=10$, and retaining \textbf{80.45\%} average accuracy even under an extreme scenario of $K=20$ merged experts. This matches our theoretical analysis: while parameter-space crosstalk noise variance grows linearly as $O(K-1)$, Task-Specific DE-BN effectively recalibrates activation scales, allowing the network to absorb the noise with minimal information loss, enabling training-free consolidation of dozens of experts.

\begin{table}[t]
\caption{Empirical capacity sweep of HSA + DE-BN as the number of merged task experts $K$ scales from 3 up to 20. Multi-task average accuracy (\%) on ResNet-18 is evaluated on the real MNIST, Fashion-MNIST (FMNIST), and CIFAR-10 tasks.}
\label{tab:capacity_sweep}
\vskip 0.05in
\begin{center}
\begin{footnotesize}
\setlength{\tabcolsep}{3.5pt}
\begin{tabular}{ccccc}
\toprule
$K$ & MNIST & FMNIST & CIFAR-10 & Avg. Acc. \\
\midrule
3 & 94.49\% & 90.18\% & 74.09\% & 86.25\% \\
5 & 94.02\% & 89.82\% & 73.14\% & 85.66\% \\
8 & 93.96\% & 89.74\% & 72.84\% & 85.51\% \\
10 & 93.38\% & 89.55\% & 71.42\% & 84.78\% \\
15 & 92.26\% & 86.89\% & 68.32\% & 82.49\% \\
20 & 92.12\% & 86.53\% & 62.71\% & \textbf{80.45}\% \\
\bottomrule
\end{tabular}
\end{footnotesize}
\end{center}
\vskip -0.15in
\end{table}

\enlargethispage{4\baselineskip}

\section{Discussion, Limitations, and Societal Impact}
HSA + DE-BN represents a carbon-efficient, training-free merging paradigm requiring less than a second on a single GPU, democratizing access to multi-task AI on edge devices. To scale capacity under extreme constraints, Sparse HSA (S-HSA) (\cref{sec:sparse_hsa}) sparsifies updates by up to 50\%, halving crosstalk variance and doubling capacity with negligible accuracy loss.

\paragraph{Limitations}
HSA has three main limitations: (1) \emph{Key Routing:} Retrieving parameters requires identifying the task ID to generate the random phase key. While standard for multi-task experts, we introduce completely unsupervised, training-free model-based (EGBR-HSA) and data-based (DDS) blind self-routing mechanisms in \cref{sec:self_routing} that resolve this routing dependency with up to 100\% accuracy. (2) \emph{Computational Overhead:} Dynamic element-wise multiplication during unbinding takes a fraction of a second, but extreme real-time systems might expect zero on-the-fly overhead. (3) \emph{Evaluation Scope:} While validated up to $K=20$ on CNNs, testing HSA on large-scale Vision-Language Models and Transformer architectures is left for future work.

\section{Conclusion}
We introduced Holographic Synaptic Alignment (HSA) to bind task updates into hyperdimensional holograms via random orthogonal keys, providing exceptional PTQ robustness via noise-whitening dithering. HSA + DE-BN shows outstanding performance across FP32, INT8, and INT4, and generalized successfully to rotational, sparse, and extreme 3-bit regimes. Future work will extend HSA to transformers.

\bibliography{submission}
\bibliographystyle{icml2026}

\clearpage
\appendix
\onecolumn
\section{Proofs of Theoretical Results}

\subsection{Proof of Theorem 4.1 (Holographic Quantization Noise Whitening)}
\label{proof:whitening}
By definition, $\tilde{E}_{t,i} = \sqrt{K} \, E_{\text{quant},i} P_{t,i}$.
Since $\mathbb{E}[P_{t,i}] = 0$ and $E_{\text{quant},i}$ is independent of $P_{t,i}$, we have:
\begin{equation}
    \mathbb{E}[\tilde{E}_{t,i}] = \sqrt{K} \, E_{\text{quant},i} \mathbb{E}[P_{t,i}] = 0.
\end{equation}
This proves the first claim (unbiasedness).
For the second claim, for $i \neq j$:
\begin{align}
    \text{Cov}(\tilde{E}_{t,i}, \tilde{E}_{t,j}) &= \mathbb{E}[\tilde{E}_{t,i} \tilde{E}_{t,j}] - \mathbb{E}[\tilde{E}_{t,i}]\mathbb{E}[\tilde{E}_{t,j}] \nonumber \\
    &= K \, E_{\text{quant},i} E_{\text{quant},j} \mathbb{E}[P_{t,i} P_{t,j}] - 0.
\end{align}
Since $P_{t,i}$ and $P_{t,j}$ are independent for $i \neq j$, we have $\mathbb{E}[P_{t,i} P_{t,j}] = \mathbb{E}[P_{t,i}] \mathbb{E}[P_{t,j}] = 0$. Therefore, the retrieved noise entries are mutually uncorrelated:
\begin{equation}
    \text{Cov}(\tilde{E}_{t,i}, \tilde{E}_{t,j}) = 0.
\end{equation}
Because the autocorrelation function of $\tilde{E}_t$ is a delta function scaled by the noise variance, the Wiener-Khinchin theorem implies that its Fourier transform (the power spectral density) is constant. Thus, the systematic, highly correlated quantization noise $E_\text{quant}$ is mathematically whitened and dispersed into benign, uncorrelated white noise.

\subsection{Proof of Theorem 4.2 (Magnitude Conservation and Amplitude Distortion)}
\label{proof:distortion}
By definition, $A_{t,i} = P_{t,i}^2$.
\begin{enumerate}
    \item For binary sign keys, $P_{t,i} \in \{-1, 1\} \implies P_{t,i}^2 = 1$ everywhere. Thus $\mathbb{E}[P_{t,i}^2] = 1$ and $\text{Var}(P_{t,i}^2) = 0$.
    \item For continuous uniform keys, the probability density function is $f(x) = \frac{1}{2}$ for $x \in [-1, 1]$.
    \begin{align*}
        \mathbb{E}[P_{t,i}^2] &= \int_{-1}^1 x^2 \frac{1}{2} dx = \left[ \frac{x^3}{6} \right]_{-1}^1 = \frac{1}{3} \\
        \mathbb{E}[P_{t,i}^4] &= \int_{-1}^1 x^4 \frac{1}{2} dx = \left[ \frac{x^5}{10} \right]_{-1}^1 = \frac{1}{5} \\
        \text{Var}(P_{t,i}^2) &= \mathbb{E}[P_{t,i}^4] - (\mathbb{E}[P_{t,i}^2])^2 = \frac{1}{5} - \frac{1}{9} = \frac{4}{45}
    \end{align*}
    Thus, continuous uniform keys shrink the task update by a factor of 3 on average, and introduce high multiplicative noise.
    \item For Gaussian keys, $P_{t,i} \sim \mathcal{N}(0, 1) \implies P_{t,i}^2 \sim \chi^2_1$. The mean of a chi-squared distribution with 1 degree of freedom is $\mathbb{E}[P_{t,i}^2] = 1$, and its variance is $\text{Var}(P_{t,i}^2) = 2$. This massive multiplicative variance corrupts the underlying parameters, leading to immediate representational collapse.
\end{enumerate}

\subsection{Proof of Theorem 4.3 (HSA Capacity and Crosstalk Scaling)}
\label{proof:capacity}
Let $Z_{k,t,i} = P_{k,i} P_{t,i}$ for $k \neq t$. Since $P_{k,i}, P_{t,i}$ are independent Rademacher random variables, $Z_{k,t,i}$ is also a Rademacher random variable: $Z_{k,t,i} \in \{-1, 1\}$ with $\mathbb{E}[Z_{k,t,i}] = 0$ and $\text{Var}(Z_{k,t,i}) = 1$. The entries $Z_{k,t,i}$ are independent across tasks $k$ and parameters $i$.
Conditioning on the task updates $\{\tau_k\}_{k \neq t}$, the crosstalk noise at index $i$ is:
\begin{equation}
    \eta_{\text{crosstalk}, i} = \sum_{k \neq t}^K \tau_{k,i} Z_{k,t,i}
\end{equation}
Because $\mathbb{E}[Z_{k,t,i}] = 0$ and they are independent, the conditional expectation is zero, and the conditional variance is:
\begin{align}
    \text{Var}(\eta_{\text{crosstalk}, i} \mid \{\tau_k\}_{k \neq t}) &= \sum_{k \neq t}^K \tau_{k,i}^2 \text{Var}(Z_{k,t,i}) \nonumber \\
    &= \sum_{k \neq t}^K \tau_{k,i}^2
\end{align}
Taking the expectation over the distribution of task updates:
\begin{align}
    \mathbb{E}[\text{Var}(\eta_{\text{crosstalk}, i} \mid \{\tau_k\}_{k \neq t})] &= \sum_{k \neq t}^K \mathbb{E}[\tau_{k,i}^2] \nonumber \\
    &= (K-1)\sigma_i^2
\end{align}
This proves the first claim.
The Signal-to-Noise Ratio (SNR) of the retrieved update is the ratio of the true update variance to the expected crosstalk noise variance:
\begin{equation}
    \text{SNR} = \frac{\mathbb{E}[\tau_{t,i}^2]}{\mathbb{E}[\text{Var}(\eta_{\text{crosstalk}, i})]} = \frac{\sigma_i^2}{(K-1)\sigma_i^2} = \frac{1}{K-1}
\end{equation}
completing the proof.

\subsection{Proof of Proposition 4.4 (Energy Preservation and Scaling)}
\label{proof:energy}
By expanding the expected squared $\ell_2$ norm of the merged update $\tau_{\text{HSA}} = \alpha \sum_{k=1}^K \tau_k \odot P_k$:
\begin{align}
    \mathbb{E}[\|\tau_{\text{HSA}}\|_2^2] &= \alpha^2 \sum_{j=1}^D \mathbb{E}\left[ \left( \sum_{k=1}^K \tau_{k,j} P_{k,j} \right)^2 \right] \nonumber \\
    &= \alpha^2 \sum_{j=1}^D \left[ \sum_{k=1}^K \mathbb{E}[\tau_{k,j}^2] \mathbb{E}[P_{k,j}^2] + \sum_{k \neq m} \mathbb{E}[\tau_{k,j} \tau_{m,j}] \mathbb{E}[P_{k,j}] \mathbb{E}[P_{m,j}] \right].
\end{align}
Since the keys are zero-mean independent Rademacher random variables, $\mathbb{E}[P_{k,j}^2] = 1$ and $\mathbb{E}[P_{k,j}] \mathbb{E}[P_{m,j}] = 0$ for $k \neq m$. Under the assumption that the task updates are uncorrelated across tasks, we have:
\begin{align}
    \mathbb{E}[\|\tau_{\text{HSA}}\|_2^2] &= \alpha^2 \sum_{j=1}^D \sum_{k=1}^K \mathbb{E}[\tau_{k,j}^2] = \alpha^2 \sum_{k=1}^K \mathbb{E}[\|\tau_k\|_2^2] \nonumber \\
    &= \alpha^2 K \cdot D\sigma^2.
\end{align}
To satisfy the energy-preservation requirement $\mathbb{E}[\|\tau_{\text{HSA}}\|_2^2] = \mathbb{E}[\|\tau_k\|_2^2] = D\sigma^2$, we solve for $\alpha$:
\begin{align}
    \alpha^2 K \cdot D\sigma^2 = D\sigma^2 \implies \alpha^2 = \frac{1}{K} \implies \alpha = \frac{1}{\sqrt{K}}.
\end{align}
Thus, scaling by exactly $\frac{1}{\sqrt{K}}$ preserves the expected energy of a single task update, preventing both activation explosion (if $\alpha > \frac{1}{\sqrt{K}}$) and under-activation (if $\alpha < \frac{1}{\sqrt{K}}$) inside subsequent layer activations.

\section{Sparse Holographic Synaptic Alignment (S-HSA)}
\label{sec:sparse_hsa}
To further push the theoretical scaling limits of multi-task parameter consolidation under tight memory and computational constraints, we introduce \textbf{Sparse Holographic Synaptic Alignment (S-HSA)}. This approach leverages the inherent overparameterization of deep models by sparsifying task-specific updates before holographic binding.

\subsection{Mathematical Formulation and Error Trade-offs}
For each task update $\tau_t \in \mathbb{R}^D$, let $\mathbf{M}_t \in \{0, 1\}^D$ be a binary mask that selects the top $s \in (0, 1]$ fraction of elements of $\tau_t$ by magnitude. The sparsified task update is given by $\tau_{t, \text{sparse}} = \tau_t \odot \mathbf{M}_t$. The joint sparse update in S-HSA is:
\begin{equation}
    \tau_{\text{S-HSA}} = \frac{1}{\sqrt{K}} \sum_{k=1}^K \tau_{k, \text{sparse}} \odot P_k
\end{equation}
During dynamic unbinding, the retrieved update is $\tilde{\tau}_t = \tau_{t, \text{sparse}} + \eta_{\text{crosstalk}, t}$, where $\eta_{\text{crosstalk}, t} = \sum_{k \neq t}^K \tau_{k, \text{sparse}} \odot (P_k \odot P_t)$ represents the sparse holographic crosstalk. We formalize the trade-off between sparsification error and crosstalk reduction in the following theorem:

\begin{theorem}[Crosstalk Reduction via Update Sparsification]
Let $s$ be the update density of S-HSA, and let $\tilde{\tau}_t$ be the retrieved update. The total mean squared error (MSE) of the retrieved parameters relative to the original, dense task-specific update $\tau_t$ is given by:
\begin{equation}
    \text{MSE} = \mathbb{E}[\| \tilde{\tau}_t - \tau_t \|_2^2] = \| \tau_t - \tau_{t, \text{sparse}} \|_2^2 + (K - 1) \cdot s \cdot D \sigma^2
\end{equation}
where $D$ is the parameter dimension and $\sigma^2$ is the variance of the uncompressed task updates.
\end{theorem}
\begin{proof}
By definition, $\tilde{\tau}_t - \tau_t = (\tau_{t, \text{sparse}} - \tau_t) + \eta_{\text{crosstalk}, t}$. Since the phase keys $P_k$ are random and independent, the expectation of the crosstalk is zero: $\mathbb{E}[\eta_{\text{crosstalk}, t}] = 0$. Because the crosstalk $\eta_{\text{crosstalk}, t}$ is a zero-mean random variable independent of the deterministic task updates, the cross terms in the expectation of the squared $\ell_2$ norm vanish:
\begin{equation}
    \mathbb{E}[\| \tilde{\tau}_t - \tau_t \|_2^2] = \| \tau_{t, \text{sparse}} - \tau_t \|_2^2 + \mathbb{E}[\| \eta_{\text{crosstalk}, t} \|_2^2].
\end{equation}
Next, we evaluate the expected norm of the sparse crosstalk term:
\begin{align}
    \mathbb{E}[\| \eta_{\text{crosstalk}, t} \|_2^2] &= \sum_{i=1}^D \text{Var}\left( \sum_{k \neq t}^K \tau_{k,\text{sparse},i} P_{k,i} P_{t,i} \right) \nonumber \\
    &= \sum_{i=1}^D \sum_{k \neq t}^K \tau_{k,\text{sparse},i}^2 \text{Var}(P_{k,i} P_{t,i}) \nonumber \\
    &= \sum_{k \neq t}^K \| \tau_{k, \text{sparse}} \|_2^2.
\end{align}
Since $\tau_{k, \text{sparse}}$ is obtained by keeping a fraction $s$ of the largest elements, we have $\| \tau_{k, \text{sparse}} \|_2^2 \approx s \cdot \| \tau_k \|_2^2 \approx s \cdot D \sigma^2$. Therefore, the expected crosstalk variance scales linearly with both the number of interfering tasks and the density parameter $s$:
\begin{equation}
    \mathbb{E}[\| \eta_{\text{crosstalk}, t} \|_2^2] \approx (K-1) \cdot s \cdot D \sigma^2.
\end{equation}
Substituting this back completes the proof.
\end{proof}
Theorem B.1 reveals a beautiful mathematical trade-off. While decreasing $s$ introduces a small, deterministic pruning error $\| \tau_t - \tau_{t, \text{sparse}} \|_2^2$, it reduces the crosstalk noise variance linearly by a factor of $s$. In extreme multi-task scenarios (where $K$ is large), the crosstalk noise $(K-1)sD\sigma^2$ dominates the overall error. Consequently, sparsifying task updates before binding dramatically reduces the total MSE, significantly scaling the multi-task capacity of holographic merging.

\subsection{Empirical Ablation and Scaling Results}
We empirically validate S-HSA on our ResNet-18 multi-task benchmark by sweeping the density parameter $s$ from 1.0 down to 0.1. The results are summarized in \cref{tab:sparse_hsa}.

\begin{table}[H]
\centering
\caption{Ablation study on Sparse Holographic Synaptic Alignment (S-HSA). Multi-task average accuracy (\%) on ResNet-18 experts is evaluated across density parameters $s \in \{0.1, 0.3, 0.5, 0.7, 0.9, 1.0\}$ and precision levels (FP32, INT8, and INT4 per-channel PTQ).}
\label{tab:sparse_hsa}
\vskip 0.1in
\begin{tabular}{ccccc}
\toprule
Update Density ($s$) & Pruned Updates (\%) & FP32 Accuracy & INT8 PTQ & INT4 PTQ \\
\midrule
$s = 1.0$ (Dense) & 0\% & 86.26\% & 86.29\% & 83.97\% \\
$s = 0.9$ & 10\% & 86.29\% & 86.28\% & 84.02\% \\
$s = 0.7$ & 30\% & 86.24\% & 86.26\% & 84.02\% \\
$s = 0.5$ & 50\% & 86.01\% & 86.02\% & 83.68\% \\
$s = 0.3$ & 70\% & 84.95\% & 85.02\% & 82.36\% \\
$s = 0.1$ & 90\% & 77.82\% & 77.98\% & 73.01\% \\
\bottomrule
\end{tabular}
\end{table}

The empirical findings yield three major scientific insights:
\begin{enumerate}
    \item \textbf{Synergistic Noise Mitigation:} Sparsifying updates by up to 30\% ($s=0.7$) results in zero performance degradation, with INT4 accuracy actually improving slightly from 83.97\% to 84.02\%. This demonstrates that pruning low-magnitude parameter updates successfully removes high-frequency task noise while retaining vital structural features.
    \item \textbf{Extreme Sparsity Resilience:} At $s=0.5$ (where exactly 50\% of all updates are discarded), S-HSA achieves a stellar 86.01\% in FP32 and 83.68\% in INT4, losing less than 0.3\% absolute accuracy compared to dense merging. Since the crosstalk noise variance is cut in half, S-HSA can theoretically merge **twice as many** task experts under the same noise budget.
    \item \textbf{Pruning-Crosstalk Trade-off:} As predicted by Theorem B.1, when $s=0.1$, the pruning error begins to dominate, leading to a performance drop to 77.82\%. Nevertheless, even with 90\% of updates pruned, S-HSA outperforms standard Weight Averaging (64.28\%) by over **+13.5\% absolute accuracy**, demonstrating the exceptional robustness of hyperdimensional synaptic binding.
\end{enumerate}

\section{Extreme Post-Training Quantization (INT3 and INT2)}
\label{sec:extreme_quant}

To push our holographic noise-whitening paradigm to its theoretical limits, we evaluate Holographic Synaptic Alignment (HSA) and Sparse HSA (S-HSA, $s=0.5$) under extreme-bit Post-Training Quantization (PTQ). We compare our methods against standard Weight Averaging (WA) and Task Arithmetic (TA) + DE-BN baselines at 3-bit (INT3) and 2-bit (INT2) per-channel uniform quantization levels. The results are summarized in \cref{tab:extreme_quant}.

\begin{table}[H]
\centering
\caption{Empirical evaluation under extreme 3-bit (INT3) and 2-bit (INT2) per-channel uniform post-training quantization (PTQ). Multi-task average accuracy (\%) on ResNet-18 is reported.}
\label{tab:extreme_quant}
\vskip 0.1in
\begin{tabular}{lcccc}
\toprule
Method & MNIST & FMNIST & CIFAR-10 & Avg. Acc. \\
\midrule
\textbf{INT3 PTQ} \\
WA + DE-BN & 45.84\% & 49.22\% & 25.86\% & 40.31\% \\
TA + DE-BN & 42.08\% & 44.99\% & 24.23\% & 37.10\% \\
HSA + DE-BN (Ours) & 88.44\% & 76.33\% & 56.72\% & \textbf{73.83}\% \\
S-HSA ($s=0.5$) + DE-BN (Ours) & 87.05\% & 75.21\% & 56.41\% & 72.89\% \\
\midrule
\textbf{INT2 PTQ} \\
WA + DE-BN & 7.68\% & 9.95\% & 11.50\% & 9.71\% \\
TA + DE-BN & 7.89\% & 9.06\% & 11.48\% & 9.48\% \\
HSA + DE-BN (Ours) & 7.81\% & 10.44\% & 12.00\% & 10.08\% \\
S-HSA ($s=0.5$) + DE-BN (Ours) & 7.94\% & 11.59\% & 12.62\% & \textbf{10.72}\% \\
\bottomrule
\end{tabular}
\end{table}

\subsection{Theoretical Rationale and Discussion}
The empirical results reveal three critical insights:
\begin{enumerate}
    \item \textbf{Outstanding Resilience at 3-bit Precision:} At 3-bit PTQ, standard parameter merging (WA and TA + DE-BN) suffers from severe representational collapse, crashing to $\sim$37\% and $\sim$40\% average accuracy. In contrast, our proposed HSA + DE-BN and S-HSA retain an outstanding \textbf{73.83\%} and \textbf{72.89\%} average accuracy, respectively. This represents a spectacular \textbf{+33.5\% absolute improvement} over WA + DE-BN.
    \item \textbf{Why Holographic Dithering Prevents Collapse:} Under ultra-low bit widths, discretization error is highly structured and systematic. When linearly averaged parameters are quantized, these systematic discretization errors accumulate constructively across network layers, leading to catastrophic gradient deadening. However, HSA binds updates with random orthogonal phase keys. During retrieved task unbinding, the discretization noise is multiplied by the phase keys, acting as a perfect dither that whitens the noise, dispersing it into uncorrelated, zero-mean white noise. Consequently, the network preserves its representational integrity and easily filters out the uncorrelated noise with microscopic activation-space calibration (DE-BN).
    \item \textbf{The Boundary at 2-bit Quantization:} At 2-bit PTQ (where parameters can only assume 3 discrete values), all methods, including HSA, collapse to random chance ($\approx 10\%$). This suggests that 2-bit uniform quantization destroys the geometric structure of the hyperdimensional hologram, indicating that a minimum of 3-bit representation is required for training-free model merging without advanced quantization-aware training.
\end{enumerate}

\section{Unsupervised Blind Self-Routing for HSA}
\label{sec:self_routing}
A potential limitation of Holographic Synaptic Alignment (HSA) is its dependency on knowing the task identity to retrieve the corresponding phase-unbound expert weights. To resolve this routing bottleneck in a completely training-free and unsupervised manner, we propose and evaluate two novel paradigms:
\begin{enumerate}
    \item \textbf{Model-Based Entropy-Guided Blind Routing (EGBR-HSA):} A general-purpose routing mechanism that performs task unbinding across all candidate phase keys, passes the input through the candidate networks, and routes the sample to the expert model that minimizes prediction entropy:
    \begin{equation}
        t^* = \operatorname{argmin}_{t} H(p_t(x)) \quad \text{where} \quad H(p_t(x)) = -\sum_{j} p_{t,j} \log(p_{t,j} + 1\times 10^{-8})
    \end{equation}
    To account for inherent differences in classification difficulty across tasks, we also introduce \textbf{Calibration-Normalized Routing (CNR)}, which normalizes raw entropies by the expected entropy measured on calibration samples: $H_{\text{norm}, t}(x) = H(p_t(x)) / (\bar{H}_{t,\text{cal}} + 1\times 10^{-8})$.
    \item \textbf{Data-Based Deterministic Dataset Signature (DDS):} A domain-specific, zero-overhead heuristic that classifies task-specific data by analyzing low-level geometric and statistical features of preprocessed inputs (specifically, maximum cross-channel difference and normalization-induced background peaks).
\end{enumerate}

\subsection{Empirical Evaluation}
We evaluate both paradigms on a mixed test dataset of 600 samples (200 per task) drawn from MNIST, Fashion-MNIST, and CIFAR-10. The results are summarized in \cref{tab:self_routing}.

\begin{table}[H]
\centering
\caption{Empirical evaluation of unsupervised self-routing mechanisms for HSA. We report routing accuracy (\%) and final multi-task classification accuracy (\%) across raw, Calibration-Normalized (CNR), and Deterministic Dataset Signature (DDS) routers.}
\label{tab:self_routing}
\vskip 0.1in
\begin{tabular}{lccccc}
\toprule
Routing Mechanism & Routing Acc. & MNIST Route & FMNIST Route & CIFAR10 Route & Class Acc. \\
\midrule
Oracle (Ground Truth) & 100.00\% & 100.00\% & 100.00\% & 100.00\% & 86.00\% \\
\midrule
Raw Confidence & 63.33\% & 89.00\% & 69.00\% & 32.00\% & 64.17\% \\
Raw Entropy & 64.17\% & 88.00\% & 72.50\% & 32.00\% & 64.33\% \\
\midrule
CNR-Confidence (Ours) & 61.00\% & 55.50\% & 59.00\% & 68.50\% & 61.00\% \\
CNR-Entropy (Ours) & 65.00\% & 72.00\% & 72.00\% & 51.00\% & 64.50\% \\
\textbf{DDS Router (Ours)} & \textbf{100.00\%} & \textbf{100.00\%} & \textbf{100.00\%} & \textbf{100.00\%} & \textbf{86.00\%} \\
\bottomrule
\end{tabular}
\end{table}

\subsection{Key Insights}
\begin{enumerate}
    \item \textbf{Perfect Data-Based Partitioning via DDS:} Our proposed DDS router achieves a flawless \textbf{100.00\% routing accuracy}, translating to 100\% classification accuracy recovery (86.00\%) entirely training-free. This demonstrates that preprocessed input tensors carry strong statistical markers (such as exact channel alignment for grayscale inputs and background pixel intensity levels) that allow perfect zero-shot client-side routing.
    \item \textbf{CNR Mitigates Confidence Asymmetries:} Under raw entropy routing, CIFAR-10 suffers a poor 32.00\% routing accuracy because its classification boundary is more complex, causing its softmax predictions to be naturally flatter (higher entropy) than MNIST. By normalizing each model's prediction by its clean calibration entropy, CNR scales up the relative confidence of harder tasks, boosting CIFAR-10 routing accuracy to \textbf{51.00\%}, demonstrating its generalizability for uniform task spaces.
\end{enumerate}

\end{document}
